% GENERATED FILE. Edit canonical lessons, then run npm run generate:books.
% canonical-source-hash: fnv1a64:3880d7fb0c3f2029
% canonical-lessons: GU-C102-sakvu, GU-C102-pase, GU-C102-takat, GU-C102-sakti, GU-C102-hak

\chapter{To be able to, With, Strength, Power, Right}
\label{ch:gu-to-be-able-to-with-strength-power-right}
\hlchaptermodality{102}

\begin{chapteropening}
\textbf{By the end of this chapter:} \emph{I can say to be able to, with, strength, power and a right.}

Complete the last lesson of chapter 102: I can say to be able to, with, strength, power and a right.
\end{chapteropening}

\section[\texorpdfstring{\gu{શકવું} (śakvũ)}{શકવું (śakvũ)}]{\gu{શકવું} (śakvũ) — to be able to, can}

\label{lesson:GU-C102-sakvu}



\begin{quote}
Before the new one: say the Gujarati for here, then the Gujarati for there.
\end{quote}

\subsection*{You'll want to know: \gu{શકવું}}
\textbf{\gu{શકવું}} (\emph{śakvũ}) — \textquotedblleft{}to be able to, can\textquotedblright{}.

\begin{grammarlens}[title={What you've built}]
One more word for this chapter.
\end{grammarlens}

\subsection*{Your turn}
\begin{itemize}
\raggedright
  \item \emph{Say it:} \emph{śakvũ}
  \item \emph{Say it:} \emph{śakvũ}, once more
  \item \emph{Say it:} say \emph{śakvũ}
  \item \emph{Recall:} say the Gujarati for here, then the Gujarati for there, then say \emph{śakvũ} again
\end{itemize}

\begin{tcolorbox}[breakable,colback=teal!4,colframe=teal!35!black,title={Before you move on}]
What does \gu{શકવું} mean? (\textquotedblleft{}To be able to, can\textquotedblright{}.) Say it once more.
\end{tcolorbox}
\section[\texorpdfstring{\gu{પાસે} (pāse)}{પાસે (pāse)}]{\gu{પાસે} (pāse) — with, in one's possession}

\label{lesson:GU-C102-pase}



\begin{quote}
Before the new one: say the Gujarati for there, then the Gujarati for to be able to.
\end{quote}

\subsection*{You'll want to know: \gu{પાસે}}
\textbf{\gu{પાસે}} (\emph{pāse}) — \textquotedblleft{}with, in one's possession\textquotedblright{}.

\begin{grammarlens}[title={What you've built}]
One more word for this chapter.
\end{grammarlens}

\subsection*{Your turn}
\begin{itemize}
\raggedright
  \item \emph{Say it:} \emph{pāse}
  \item \emph{Say it:} \emph{pāse}, once more
  \item \emph{Say it:} say \emph{pāse}
  \item \emph{Recall:} say the Gujarati for there, then the Gujarati for to be able to, then say \emph{pāse} again
\end{itemize}

\begin{tcolorbox}[breakable,colback=teal!4,colframe=teal!35!black,title={Before you move on}]
What does \gu{પાસે} mean? (\textquotedblleft{}With, in one's possession\textquotedblright{}.) Say it once more.
\end{tcolorbox}
\section[\texorpdfstring{\gu{તાકાત} (tākāt)}{તાકાત (tākāt)}]{\gu{તાકાત} (tākāt) — strength}

\label{lesson:GU-C102-takat}



\begin{quote}
Before the new one: say the Gujarati for to be able to, then the Gujarati for with.
\end{quote}

\subsection*{You'll want to know: \gu{તાકાત}}
\textbf{\gu{તાકાત}} (\emph{tākāt}) — \textquotedblleft{}strength\textquotedblright{}.

\begin{grammarlens}[title={What you've built}]
One more word for this chapter.
\end{grammarlens}

\subsection*{Your turn}
\begin{itemize}
\raggedright
  \item \emph{Say it:} \emph{tākāt}
  \item \emph{Say it:} \emph{tākāt}, once more
  \item \emph{Say it:} say \emph{tākāt}
  \item \emph{Recall:} say the Gujarati for to be able to, then the Gujarati for with, then say \emph{tākāt} again
\end{itemize}

\begin{tcolorbox}[breakable,colback=teal!4,colframe=teal!35!black,title={Before you move on}]
What does \gu{તાકાત} mean? (\textquotedblleft{}Strength\textquotedblright{}.) Say it once more.
\end{tcolorbox}
\section[\texorpdfstring{\gu{શક્તિ} (śakti)}{શક્તિ (śakti)}]{\gu{શક્તિ} (śakti) — power}

\label{lesson:GU-C102-sakti}



\begin{quote}
Before the new one: say the Gujarati for with, then the Gujarati for strength.
\end{quote}

\subsection*{You'll want to know: \gu{શક્તિ}}
\textbf{\gu{શક્તિ}} (\emph{śakti}) — \textquotedblleft{}power\textquotedblright{}.

\begin{grammarlens}[title={What you've built}]
One more word for this chapter.
\end{grammarlens}

\subsection*{Your turn}
\begin{itemize}
\raggedright
  \item \emph{Say it:} \emph{śakti}
  \item \emph{Say it:} \emph{śakti}, once more
  \item \emph{Say it:} say \emph{śakti}
  \item \emph{Recall:} say the Gujarati for with, then the Gujarati for strength, then say \emph{śakti} again
\end{itemize}

\begin{tcolorbox}[breakable,colback=teal!4,colframe=teal!35!black,title={Before you move on}]
What does \gu{શક્તિ} mean? (\textquotedblleft{}Power\textquotedblright{}.) Say it once more.
\end{tcolorbox}
\section[\texorpdfstring{\gu{હક} (hak)}{હક (hak)}]{\gu{હક} (hak) — a right}

\label{lesson:GU-C102-hak}



\begin{quote}
Before the new one: say the Gujarati for strength, then the Gujarati for power.
\end{quote}

\subsection*{You'll want to know: \gu{હક}}
\textbf{\gu{હક}} (\emph{hak}) — \textquotedblleft{}a right\textquotedblright{}.

\begin{grammarlens}[title={What you've built}]
One more word for this chapter.
\end{grammarlens}

\subsection*{Your turn}
\begin{itemize}
\raggedright
  \item \emph{Say it:} \emph{hak}
  \item \emph{Say it:} \emph{hak}, once more
  \item \emph{Say it:} say \emph{hak}
  \item \emph{Recall:} say the Gujarati for strength, then the Gujarati for power, then say \emph{hak} again
\end{itemize}

\begin{tcolorbox}[breakable,colback=teal!4,colframe=teal!35!black,title={Before you move on}]
What does \gu{હક} mean? (\textquotedblleft{}A right\textquotedblright{}.) Say it once more.
\end{tcolorbox}
